% ===================================================================
%  도표 제 5 호 — 집과 그 짓기.
%
%  Traduction, faite en 2026, en coreen d'aujourd'hui, sur l'ido de
%  texto/io. Regles de la colonne : voir 10-dopyo-01.tex. Le tableau
%  entier est un DIALOGUE et aucun alinea n'est numerote par
%  l'auteur : toutes les cles portent « 01 », le rang suivant la
%  replique, comme en ido.
%
%  LE TABLEAU LE PLUS LONG DU LIVRET — 68 blocs, plus de cent vingt
%  renvois — ET LE PREMIER QUE J'ECRIS AVEC LA LISTE DE parigi.py EN
%  MAIN DES LA PREMIERE LIGNE. C'est exactement ce qui avait ete fait
%  pour la colonne telougoue au meme tableau, et pour la meme raison :
%  c'est ici que l'ordre des renvois coute le plus cher si on le
%  decouvre apres coup.
%
%  DOUZE RETOURNEMENTS SUR TRENTE-SEPT PAIRES, et le tri des
%  vingt-cinq autres est celui qu'on connait : enumeration
%  (« sablo (62), kalko (63), aquo (64) »), frontiere de phrase, et
%  renvoi de gauche deja modificateur. Douze sur trente-sept est le
%  taux le plus BAS de la colonne coreenne jusqu'ici — plus bas que
%  le defile du tableau 3 (dix-sept sur vingt-cinq) — et la raison
%  est la meme qu'au marche telougou : un chantier s'enumere. On
%  nomme les corps de metier l'un apres l'autre, et les outils de
%  chacun a la suite.
%
%  L'ARTISAN COREEN A SON NOM, ET IL EST VIEUX. C'est le tableau ou
%  la langue se montre le plus entiere : 목수 le charpentier, 미장이
%  le platrier, 석공 le macon, 대장장이 le forgeron, 소목장이 le
%  menuisier, 자물쇠장이 le serrurier, 유리장이 le vitrier, 도배장이
%  le tapissier, 칠장이 le peintre, 정원사 le jardinier, 마부 le
%  garcon d'ecurie. Les suffixes -장이 et -꾼 font a eux seuls la
%  moitie de la liste, et aucun n'est un calque.
%
%  Les outils suivent : 흙손 la truelle, 다림추 le fil a plomb, 대패
%  le rabot, 끌 le ciseau, 모루 l'enclume, 집게 les tenailles, 갈퀴
%  le rateau, 물뿌리개 l'arrosoir, 빗장 le verrou, 홈통 la gouttiere,
%  피뢰침 le paratonnerre, 바람개비 la girouette, 박공 le pignon.
%  Trois emprunts seulement dans tout le tableau, et tous pour des
%  materiaux : 슬레이트 l'ardoise, 페인트 la peinture, 퍼티 le mastic.
%
%  DEUX RENVOIS QUE LE FAC-SIMILE DONNE SANS GRAS — (58) la pompe,
%  (69) le bois, (125) le forgeron, (34) le grenier — restent sans
%  gras ici, comme dans les trente-huit colonnes precedentes. Et le
%  (107) est donne DEUX FOIS, au porteur de mortier puis au
%  tapissier : faute de l'edition ido, que la source ne corrige pas.
% ===================================================================

%%K t05-apar-1 apar
\VUsaut{6.0mm}
\VUtitre{15.0pt}{60}{도표 제 5 호}
\VUsaut{1.4mm}
\VUtitre{11.4pt}{0}{집과 그 짓기.}
\VUsaut{2.2mm}

%%K t05-tit-1 sub
\VUpk{10.2pt}{40}{반가운 만남.}
\VUsaut{3.0mm}

%%K t05-01-1 p
\textsc{요안네스}. --- 아저씨, 저는 \VUgras{집}을 어떻게 짓는지 몹시 알고 싶어요.

%%K t05-01-2 p
\textsc{아저씨} \textsuperscript{(80)}. --- 그건 아주 쉽단다, 얘야. 나와 함께 가자.
베르나르두스 씨 \textsuperscript{(79)}가 짓게 하고 있고 내가 살게 될 집을 보여 주마.

%%K t05-01-3 p
\textsc{요안네스}. --- 그런데 이 집의 \VUgras{도면} \textsuperscript{(76)}은 누가
그렸어요?

%%K t05-01-4 p
\textsc{아저씨}. --- \VUgras{건축가} \textsuperscript{(75)}란다. 그가 아주 또렷한 그림을
내놓아 그것이 받아들여졌고, \VUgras{도급업자} \textsuperscript{(77)}가 일을 시작했지.

%%K t05-01-5 p
\textsc{요안네스}. --- 그럼 도급업자는 누구예요?

%%K t05-01-6 p
\textsc{아저씨}. --- 블랑코 씨란다. 네 동무 야코부스의 아버지시지. 그분이 건축가 루포
씨와 함께 일꾼들을 이끄신다.

%%K t05-01-7 p
아! 마침 저기 블랑코 씨가 자기 \VUgras{자} \textsuperscript{(78)}를 들고 오시고, 루포
씨도 도면을 들고 오시는구나.

%%K t05-01-8 p
안녕하십니까, 두 분. 어떻게 지내십니까?

%%K t05-01-9 p
\textsc{루포 씨}. --- 아주 잘 지냅니다, 고맙습니다. 안녕, 요안네스. 이 아이가 참 많이
컸군요!

%%K t05-01-10 p
\textsc{아저씨}. --- 예, 정말로 작은 어른이 다 됐지요. 아주 진지해서, 보는 것마다 다
알려 주어야 합니다. 오늘은 집을 어떻게 짓는지 알고 싶어 하는군요.

%%K t05-tit-2 sub
\VUpk{10.2pt}{40}{일꾼들.}
\VUsaut{3.0mm}

%%K t05-01-11 p
\textsc{블랑코 씨}. --- 자, 나와 함께 가자, 얘야. 곧 설명해 주마. 나는 배우려 드는
아이들을 아주 좋아한단다.

%%K t05-01-12 p
손에 \VUgras{삽} \textsuperscript{(82)}을 든 저 일꾼이 보이지. 저 사람이
\VUgras{땅파는 일꾼} \textsuperscript{(81)}이다. 그는 물관을 놓으려고 \VUgras{도랑}
\textsuperscript{(46)}을 판다. 또 집이 들어설 자리의 땅도 파냈다. 석공이 네 \VUgras{벽}을
세울 수 있게 말이다. 그 벽에서 땅속에 들어가는 부분을 \VUgras{기초} \textsuperscript{(2)}
라고 한다. 기초 사이에 들어 있는 공간이 \VUgras{지하실} \textsuperscript{(3)}을 이룬다. 그
지하실에는 \VUgras{지하실 창} \textsuperscript{(5)}이라는 작은 창과 \VUgras{문}
\textsuperscript{(4)}과 층계가 있다.

%%K t05-01-13 p
\VUgras{석공} \textsuperscript{(108)}은 \VUgras{문} \textsuperscript{(8)}과 \VUgras{창}
\textsuperscript{(17)} 자리를 남겨 두고 벽을 계속 쌓았다.

%%K t05-01-14 p
\textsc{요안네스}. --- 그런데 그 석공이 안 보이는데요!

%%K t05-01-15 p
\textsc{블랑코} 씨. --- 지금은 집 안 일을 마쳤단다. 그는 \VUgras{마구간}
\textsuperscript{(54)} 가까이 제 \VUgras{비계} \textsuperscript{(110)} 위에 있다. 거기서
\VUgras{빨래터} \textsuperscript{(56)}의 벽을 쌓고 있지.

%%K t05-01-16 p
그에게는 흙손과 흙받기와 \VUgras{다림추} \textsuperscript{(109)}가 있다. 그런데 너는 그
이름들을 다 외우지 못할 게다.

%%K t05-01-17 p
\textsc{요안네스}. --- 아니에요, 외울 거예요. 곧 아저씨께 작은 흙손과 흙받기를 사
달라고 조를 거고요, 다림추는 실로 제가 만들 거예요.

%%K t05-tit-3 sub
\VUsaut{3.0mm}
\VUpk{10.2pt}{40}{재료.}
\VUsaut{3.0mm}

%%K t05-01-18 p
\textsc{아저씨}. --- 하! 아주 좋구나. 그런데 요안네스야, 집을 지으려면 무엇이 필요하냐?

%%K t05-01-19 p
\textsc{요안네스}. --- \VUgras{다듬돌} \textsuperscript{(59)}과 \VUgras{회반죽}
\textsuperscript{(65)}이 필요해요.

%%K t05-01-20 p
\textsc{아저씨}. --- 맞다. 큰 모자를 쓰고 제 \VUgras{짐수레} \textsuperscript{(136)} 위에
앉은 저 사람을 보아라. 저 사람이 \VUgras{수레꾼} \textsuperscript{(135)}이다. 날마다
여기로 \VUgras{돌덩이} \textsuperscript{(60)}를 실어 온다. 그는 수레를 마당에서 부리고,
\VUgras{석재 다듬는 이} \textsuperscript{(83)}들이 그 덩이를 다듬고 제 \VUgras{돌톱}
\textsuperscript{(85)}으로 켠다. \VUgras{막일꾼} \textsuperscript{(146)}이 그것을 날라
석공에게 주면 석공이 자리에 앉힌다.

%%K t05-01-21 p
그동안 \VUgras{반죽꾼} \textsuperscript{(87)}이 회반죽을 갠다. 그에게는 \VUgras{모래}
\textsuperscript{(62)}, \VUgras{석회} \textsuperscript{(63)}, \VUgras{물}
\textsuperscript{(64)}이 필요하다. 석회는 그 곁 \VUgras{자루} \textsuperscript{(71)}에
있고, \VUgras{석공 조수} \textsuperscript{(111)}가 \VUgras{손수레} \textsuperscript{(112)}로
모래를 날라다 준다. \VUgras{견습공} \textsuperscript{(113)}은 펌프 (58) 곁에 있다. 그가
\VUgras{물통} \textsuperscript{(114)}을 채운다. 회반죽은 곧 준비된다.
\VUgras{회반죽 나르는 이} \textsuperscript{(107)}가 그것을 받아 사다리로 비계 위까지
지고 올라가고, 거기에서 석공이 집의 그쪽 면을 마무리한다.

%%K t05-01-22 p
그러면 \VUgras{지붕} \textsuperscript{(31)} 일을 하는 일꾼은 무엇이라고 하느냐?

%%K t05-01-23 p
\textsc{요안네스}. --- \VUgras{지붕 이는 이} \textsuperscript{(118)}예요.

%%K t05-tit-4 sub
\VUsaut{3.0mm}
\VUpk{10.2pt}{40}{집의 지붕.}
\VUsaut{3.0mm}

%%K t05-01-24 p
\textsc{블랑코} 씨. --- 그렇지. 그러나 먼저 \VUgras{목수} \textsuperscript{(115)}가 온다.

%%K t05-01-25 p
목수는 \VUgras{지붕틀} \textsuperscript{(35)}을 짠다. 그는 \VUgras{들보}
\textsuperscript{(37)}를 맞춘다 — \VUgras{나무못} \textsuperscript{(117)}으로. 저 위에서
넓은 벨벳 반바지를 입고 있는 저 일꾼들이 목수다. 하나는 작은 들보를 잡고, 다른 하나는
제 \VUgras{도끼} \textsuperscript{(116)}로 나무못을 깎는다. \VUgras{굴뚝}
\textsuperscript{(41)} 가까이에 있는 사람이 지붕 이는 이다. 그는

%%K t05-noto-1 noto
\VUnotes{143.2mm}{%
(*) \VUgras{Balk-o} = 독일어 \textit{Balken}, 영어 \textit{joist}, 프랑스어
\textit{solive}, 이탈리아어 \textit{travicello}, 러시아어 \textit{balka},
스페인어와 포르투갈어 \textit{viga}.
지금까지 받아들여지지 않은 기술 어근이지만, 프랑스어 \textit{solive}의 뜻을
옮기기 위해 제안할 만하다. 그것은 프랑스어 \textit{poutre}인 큰 들보도 아니고
작은 들보도 아니다. 마루청과 천장을 받치는 네모진 나무 막대다.%
}

%%K t05-01-25 p suite
(*) 장선 위에 산자를 대고, 산자 위에 \VUgras{슬레이트} \textsuperscript{(66)}를 얹는다.
때로는 기와를 얹기도 한다.

%%K t05-01-26 p
마구간 지붕은 \VUgras{기와지붕} \textsuperscript{(55)}이고, 집 지붕은
\VUgras{슬레이트 지붕} \textsuperscript{(32)}이며, 베란다 지붕은 \VUgras{아연 지붕}
\textsuperscript{(24)}이다.

%%K t05-01-27 p
\textsc{요안네스}. --- 지붕 이는 이가 아연 지붕도 하나요?

%%K t05-01-28 p
\textsc{블랑코} 씨. --- 아니다. 그건 \VUgras{아연공} \textsuperscript{(121)}이나
\VUgras{납공}이 한다. 집 지붕에 \VUgras{지붕창} \textsuperscript{(38)}을 다는 사람이지.
그는 \VUgras{빗물 홈통} \textsuperscript{(43)}과 \VUgras{처마 홈통}
\textsuperscript{(42)}도 단다. 아연과 함석을 땜질해 붙인다. \VUgras{피뢰침}
\textsuperscript{(40)}을 붙인 것도 그이고, \VUgras{바람개비} \textsuperscript{(39)}를
붙인 것도 그다 — \VUgras{박공} \textsuperscript{(36)} 위에.

%%K t05-01-29 p
\textsc{요안네스}. --- 그러면 집이 다 된 건가요?

%%K t05-tit-5 sub
\VUsaut{3.0mm}
\VUpk{10.2pt}{40}{연장.}
\VUsaut{3.0mm}

%%K t05-01-30 p
\textsc{블랑코} 씨. --- 아직 다는 아니다. \VUgras{미장이} \textsuperscript{(99)}가
\VUgras{벽돌} \textsuperscript{(61)}로 집을 여러 방으로 나누는 칸막이를 세운다. 그는
\VUgras{회벽} \textsuperscript{(70)}으로 \VUgras{천장} \textsuperscript{(26)}과 벽을
바른다. 지금은 \VUgras{베란다} \textsuperscript{(33)}의 천장을 바르고 있다. 그에게도
석공처럼 \VUgras{흙손} \textsuperscript{(100)}과 \VUgras{흙받기} \textsuperscript{(101)}가
있다.

%%K t05-01-31 p
그 다음에 소목장이가 문과 창과 \VUgras{마루} \textsuperscript{(16)}와 \VUgras{덧문}
\textsuperscript{(19)}을 짠다.

%%K t05-01-32 p
지금 그는 마당에서 제 \VUgras{작업대} \textsuperscript{(90)} 위에서 일하고 있다.
그에게는 \VUgras{톱} \textsuperscript{(94)}이 있고 — 나무 (69)를 켜는 것이다 —,
\VUgras{대패} \textsuperscript{(91)}, \VUgras{죔쇠} \textsuperscript{(92)}, \VUgras{끌}
\textsuperscript{(94 bis)}, \VUgras{컴퍼스} \textsuperscript{(95)}, 그리고 나무
\VUgras{망치} \textsuperscript{(95 bis)}가 있다.

%%K t05-01-33 p
\textsc{요안네스}. --- 아! 그런데 소목 일이 미장 일보다 더 재미있겠어요. 아저씨, 저에게
작업대와 톱과 대패를 사 주세요. 곧 메도르에게 줄 개집을 만들 거예요.

%%K t05-01-34 p
\textsc{아저씨}. --- 그러마, 얘야.

%%K t05-01-35 p
\textsc{요안네스}. --- 정말 좋아요! 그런데 대문 곁에 서 있는 저 사람은 누구예요?

%%K t05-tit-6 sub
\VUsaut{3.0mm}
\VUpk{10.2pt}{40}{칠하기.}
\VUsaut{3.0mm}

%%K t05-01-36 p
\textsc{블랑코} 씨. --- 저 사람은 \VUgras{칠장이} \textsuperscript{(102)}다. 그는 제
사다리 위에 서서 \VUgras{페인트} \textsuperscript{(103)} 통과 제 \VUgras{붓}
\textsuperscript{(104)}을 들고 있다. 더 오래 가라고 대문의 나무를 칠하는 것이다.

%%K t05-01-37 p
그는 방마다 종이를 바르는 일도 한다.

%%K t05-01-38 p
\VUgras{도배장이} \textsuperscript{(107)}는 \VUgras{사다리} \textsuperscript{(106)} 위에
있는데, 그 사다리는 \VUgras{발코니} \textsuperscript{(28)} 위에 있다. 그는 \VUgras{차양}
\textsuperscript{(29)}을 단다.

%%K t05-01-39 p
큰 창 곁에 서 있는 일꾼은 \VUgras{유리장이} \textsuperscript{(96)}다. 그는
\VUgras{유리판} \textsuperscript{(18)}을 끼운다 — \VUgras{퍼티} \textsuperscript{(73)}로.
지하실 문 곁에 무릎을 꿇고 있는 저 다른 일꾼은 \VUgras{자물쇠장이}
\textsuperscript{(97)}다. 그는 \VUgras{자물쇠} \textsuperscript{(6)}를 단다. 그의
\VUgras{줄} \textsuperscript{(98)}이 그 곁에 있다.

%%K t05-01-40 p
\textsc{요안네스}. --- 제 \VUgras{망치} \textsuperscript{(84)}로 쇠를 두드리는 저 일꾼도
자물쇠장이인가요?

%%K t05-01-41 p
\textsc{블랑코} 씨. --- 더 바로 말하면 그는 대장장이 (125)다. 그는 \VUgras{쇠붙이}
\textsuperscript{(67)}를 벼린다 — \VUgras{전기공} \textsuperscript{(124)}에게 줄 것이다.
그 전기공은 \VUgras{차고} \textsuperscript{(51)} 지붕 위에 있다. 대장장이에게는
\VUgras{화덕} \textsuperscript{(126)}, \VUgras{모루} \textsuperscript{(127)},
\VUgras{집게} \textsuperscript{(128)}가 있다.

%%K t05-01-42 p
전기공은 전화의 \VUgras{구리} \textsuperscript{(44)} \VUgras{줄}을 맨다.

%%K t05-tit-7 sub
\VUpk{10.2pt}{40}{뜰 가꾸기.}
\VUsaut{3.0mm}

%%K t05-01-43 p
\textsc{아저씨}. --- \VUgras{마당} \textsuperscript{(45)} 안쪽, \VUgras{쇠울}
\textsuperscript{(47)}과 \VUgras{대문} \textsuperscript{(48)} 가까이에 \VUgras{정원사}
\textsuperscript{(129)}가 보이지. 그는 \VUgras{작은 뜰} \textsuperscript{(49)}을 손질하고
있다. 제 \VUgras{삽} \textsuperscript{(131)}으로 땅을 파 뒤집었고, 씨를 뿌리고
\VUgras{갈퀴} \textsuperscript{(130)}로 고른다. 그 다음에 제 \VUgras{물뿌리개}
\textsuperscript{(132)}로 \VUgras{화단} \textsuperscript{(150)}에 물을 줄 것이고, 그러면
싹이 자랄 것이다.

%%K t05-01-44 p
\VUgras{마부} \textsuperscript{(133)}는 방금 말을 돌보고 먹이를 준 뒤, \VUgras{마차}
\textsuperscript{(52)}를 닦는다 — 스펀지와 \VUgras{손펌프} \textsuperscript{(134)}와
물통으로. 그러고 나서 마차를 차고에 다시 넣고 굵은 \VUgras{빗장} \textsuperscript{(53)}
으로 문을 잠글 것이다.

%%K t05-01-45 p
자, 이만하면 넉넉한 설명을 들은 것 같구나.

%%K t05-01-46 p
\textsc{요안네스}. --- 예, 아저씨. 그래도 집 안을 보고 싶어요.

%%K t05-01-47 p
\textsc{아저씨}. --- 그건 내일 하자. 루포 씨가 도면을 설명해 주고 방마다 이름을 일러
줄 거다.

%%K t05-tit-8 sub
\VUsaut{3.0mm}
\VUpk{10.2pt}{40}{집 안의 짜임.}
\VUsaut{2.4mm}

%%K t05-apar-2 apar
{\centering\textit{(도면을 보라.)}\par}
\VUsaut{2.4mm}

%%K t05-01-48 p
\textsc{건축가}. --- 오너라, 요안네스. 곧 집의 여러 부분을 둘러보게 해 주마.

%%K t05-01-49 p
\VUgras{아래층} \textsuperscript{(7)}으로 들어가자. 여기가 \VUgras{초인종}
\textsuperscript{(11)} 단추다. \VUgras{계단} \textsuperscript{(14)} 셋을 오르는데, 그것은
\VUgras{현관 층계} \textsuperscript{(9)}의 계단이다. \VUgras{문지방}
\textsuperscript{(10)}을 넘는데, 그것은 \VUgras{대문} \textsuperscript{(8)}의 문지방이다.
그러면 \VUgras{층계} \textsuperscript{(13)} 밑에 이르게 된다.

%%K t05-01-50 p
\textsc{요안네스}. --- 여기가 복도군요.

%%K t05-01-51 p
\textsc{건축가}. --- 아니다. 여기는 \VUgras{현관방} \textsuperscript{(12)}이다.
\VUgras{복도} \textsuperscript{(a)}는 저 끝에 있다. 오른쪽에 \VUgras{응접실}
\textsuperscript{(b)}과 \VUgras{식당} \textsuperscript{(c)}이 있고, 식당 맞은편에
\VUgras{부엌} \textsuperscript{(d)}과 \VUgras{하인방} \textsuperscript{(e)}이 있으며, 그
앞쪽에 \VUgras{서재} \textsuperscript{(f)}가 있다. \VUgras{베란다} \textsuperscript{(23)}
는 응접실 곁에 있는데, \VUgras{유리문} \textsuperscript{(25)}이 달려 있다.

%%K t05-01-52 p
이제 층계를 오르자. \VUgras{난간} \textsuperscript{(15)}을 잡고 계단을 세어 보아라. 열넷
이다. 여기가 \VUgras{층계참} \textsuperscript{(m)}이고, 우리는 이제 이 \VUgras{층}
\textsuperscript{(27)}에 올라왔다.

%%K t05-tit-9 sub
\VUsaut{3.0mm}
\VUpk{10.2pt}{40}{이 층.}
\VUsaut{3.0mm}

%%K t05-01-53 p
층계 맞은편에 \VUgras{뒷간} \textsuperscript{(20)}과 \VUgras{세면실} \textsuperscript{(j)}
이 있다. 오른쪽에는 고운 \VUgras{침실} \textsuperscript{(g)}이 있는데, 집 \VUgras{정면}
\textsuperscript{(1)}으로 창이 둘 나 있다. 왼쪽에는 \VUgras{목욕실} \textsuperscript{(1)}과
\VUgras{손님방} \textsuperscript{(i)}이 있는데, 그 방에는 큰 \VUgras{벽감}
\textsuperscript{(h)}이 있다. 침실 곁에는 \VUgras{아이방} \textsuperscript{(k)}이 있다.
그 방은 넓은 \VUgras{테라스} \textsuperscript{(21)}에 면해 있다. 이 테라스 아래에 또
다른 뒷간 (20)이 있고, 벽에는 저녁 시간을 알리는 \VUgras{종} \textsuperscript{(22)}이
있다.

%%K t05-01-54 p
\textsc{요안네스}. --- \VUgras{삼 층}으로 올라가요.

%%K t05-01-55 p
\textsc{건축가}. --- 삼 층은 없다. 지붕 밑에는 \VUgras{다락방} \textsuperscript{(33)}과
곳간 (34)만 있다. 둘러보기는 끝났으니 다시 내려가도 되겠다.

%%K t05-01-56 p
\textsc{요안네스}. --- 고맙습니다, 선생님. 참으로 재미있었어요. 곧 아버지께 제가 본 것을
모두 이야기해 드릴 거예요.

\VUsaut{4.0mm}
\VUfilet{22mm}
